% C.3be. Same-ID semilocal refinement of the certified henselian version.
\paragraph{C.3be. Rational top Chow vanishing with residue field
\(\overline{\mathbb F}_p\) in every local dimension (Proved).}
Let \((R,\mathfrak m,k)\) be an excellent local domain
with \(k\simeq\overline{\mathbb F}_p\) and \(N=\dim R\geq2\).
Put \(S=\operatorname{Spec}R\) and
\(X=S\setminus\{\mathfrak m\}\). Then
\begin{equation}\label{eq:C3be-top-vanishing}
              \mathrm{CH}_1(S)_{\mathbb Q}=0,
              \qquad \mathrm{CH}^{N-1}(X)_{\mathbb Q}=0.
\end{equation}
The dimension-indexed group uses
\(\delta(\mathfrak p)=\dim(R/\mathfrak p)\), and the puncture
inherits this dimension function. Thus
\(\mathrm{CH}_1(X)=\mathrm{CH}^{N-1}(X)\); we do not reindex
using the smaller Krull dimension of the puncture.
The ring need not be henselian, normal, algebraizable over a field, or of
equal characteristic, and its puncture need not be regular.
In the setting where the ordinary algebraic-equivalence quotient
\(B^{N-1}(X)\) is defined, that quotient also vanishes rationally.
In particular this gives the ordinary top quotient for the
strict-local \(\overline{\mathbb F}_p\)-algebra germs in the
manuscript. We impose no coefficient-field family convention on
a mixed-characteristic ring.

\begin{proof}
\emph{1. A projective normal-crossings resolution of a local surface.}
We first take a normal excellent local domain
\((B,\mathfrak n,k)\) of dimension two, with the same residue
field. Let \(U=\operatorname{Spec}B\setminus\{\mathfrak n\}\).
Normality makes \(U\) regular. Lipman's surface-resolution theorem
in \citet[Theorem 54.14.5, Tag 0BGP]{frC3beStacksLipman}
applies: the normalization is finite, the only possible singular
point is the closed point, and excellence makes its completed
normal local ring normal. Its proof, together with
\citet[Lemma 54.6.5, Tag 0BFZ]{frC3beStacksNormalizedPointBlowups},
produces a finite sequence of normalized blowups at closed points
over \(\mathfrak n\), ending in a regular surface. All normalization
maps are finite by excellence. If \(B\) is regular, we start
instead with its ordinary blowup at \(\mathfrak n\).

Apply embedded resolution to the entire reduced closed fibre.
\citet[Lemma 54.15.6, Tag 0BIC]{frC3beStacksSurfaceSNC}
uses further blowups at closed points on that fibre and makes its
reduced inverse image a strict normal-crossings divisor. The
curve-normalization finiteness needed there follows from
excellence. We obtain a regular integral projective surface
\[
 f:Z\longrightarrow\operatorname{Spec}B,\qquad
 E=(f^{-1}(\mathfrak n))_{\mathrm{red}}=\bigcup_{i=1}^{c}E_i,
\]
with \(f\) an isomorphism over \(U\). Each \(E_i\) is a smooth
projective integral curve over the perfect field \(k\).

We also retain an effective exceptional Cartier divisor
\begin{equation}\label{eq:C3be-negative-ample}
 D=\sum_{i=1}^{c}a_iE_i,\qquad a_i>0,\qquad
                    \mathcal O_Z(-D)\text{ is }f\text{-ample}.
\end{equation}
Here is a direct construction, so projectivity and preservation
of \(U\) are simultaneous. On each ordinary blowup the exceptional
divisor \(F\) is effective Cartier and
\(\mathcal O(-F)=\mathcal O(1)\) is relatively ample
\citep[Section 31.33 and Lemma 31.33.4, Tag 02OS]
{frC3beStacksBlowupTautological}. Finite normalization preserves
this ampleness by pullback
\citep[Lemma 30.17.2, Tag 0B5V]{frC3beStacksFiniteAmple};
the pullback divisor remains effective Cartier because the
normalization is dominant and integral. If \(h\) is a subsequent
step and \(-D_0\) is ample for the previous composite, then
\(-F_h-a h^*D_0\) is ample for the new composite when \(a\)
is sufficiently large, by
\citet[Lemma 29.38.8, Tag 0C4K]{frC3beStacksCompositeAmple}.
The first blowup has the whole closed fibre as its exceptional
support, and its total transform retains that support at every
later step. The resulting divisor has every coefficient positive,
as required in \eqref{eq:C3be-negative-ample}. Properness and
the resulting ample invertible sheaf over the affine base also
give projectivity.

\emph{2. Finite detection of \(K_0\) and the mixed-characteristic point.}
Set \(E_b=Z\times_B\operatorname{Spec}(B/\mathfrak n^b)\).
Apply pro-descent to the abstract blowup square with centre
\(\operatorname{Spec}k\). Theorem A and exact sequence (1.6),
author pages 2--3, of \citet{rootKerzStrunkTamme2018} give
\[
 K_0(B)\longrightarrow
       \{K_0(B/\mathfrak n^b)\}_b\oplus K_0(Z)
       \longrightarrow\{K_0(E_b)\}_b.
\]
Both local-ring groups on the left are canonically \(\mathbb Z\),
and the first arrow is the diagonal graph of the structure-sheaf
class, up to the harmless sign in the exact sequence.
Its cokernel is the constant pro-group \(K_0(Z)\). Consequently
\[
               K_0(Z)\longrightarrow\{K_0(E_b)\}_b
\]
is a monomorphism in the category of pro-abelian groups. This
gives an actual finite detecting level: the levelwise kernels
are subgroups of the fixed group \(K_0(Z)\), their transition
maps are inclusions, and their pro-system is zero. Hence for
some \(b_0\) the map \(K_0(Z)\to K_0(E_{b_0})\) is injective.

The immersion \(E\hookrightarrow E_b\) is nilpotent. Moreover
\(p\in\mathfrak n\), so \(p^b=0\) on \(E_b\), even if \(B\)
has mixed characteristic. \citet[Theorem A, p. 29]
{rootGeisserHesselholt2011} assumes that \(p\) is nilpotent in
the ring; it does not require an \(\mathbb F_p\)-algebra.
Its relative groups are \(p\)-primary torsion in every degree.
Applying the theorem on a finite affine cover and using
Zariski descent \citep[Theorem 8.1, p. 367]
{D_ThomasonTrobaugh1990} therefore gives
\(K_0(E_b)_{\mathbb Q}\simeq K_0(E)_{\mathbb Q}\).
We conclude that restriction is injective:
\begin{equation}\label{eq:C3be-K0-injective}
                         K_0(Z)_{\mathbb Q}
                         \lhook\joinrel\longrightarrow
                         K_0(E)_{\mathbb Q}.
\end{equation}
No nilpotence of \(p\) on \(Z\) or on \(B\) is used.

\emph{3. Multidegree zero implies torsion for line bundles on \(Z\).}
A line bundle on \(E\) whose degree is zero on every \(E_i\)
is torsion. Indeed, the curves and any individual degree-zero
line bundle descend to a finite field, where the group of points
of their Jacobians is finite. Thus a common power trivializes
all component restrictions. After these trivializations the
remaining gluing data consist of finitely many scalars at the
nodes. Every element of \(k^*\) is torsion, so a further common
power kills these scalars as well. This also proves the assertion
if the dual graph is disconnected.

Let \(L\in\operatorname{Pic}(Z)\) have degree zero on each
\(E_i\). Choose \(M>0\) with
\(L^{\otimes M}|_E\simeq\mathcal O_E\).
By \eqref{eq:C3be-K0-injective},
\([L^{\otimes M}]-[\mathcal O_Z]\) is zero in
\(K_0(Z)_{\mathbb Q}\), and hence is killed by an integer
\(n>0\) in \(K_0(Z)\). Since \(Z\) is projective over an
affine scheme, it has the resolution property. The comparison
\(K_0(\operatorname{Vect}Z)\simeq K_0(Z)\) is
\citet[Lemma 36.38.5, Tag 0FDJ]{frC3beStacksVectorK0}.
The determinant homomorphism
\citep[Section 17.26 and Lemma 17.26.1, Tag 0B37]
{frC3beStacksDeterminant} gives
\(L^{\otimes Mn}\simeq\mathcal O_Z\). We have proved that
the multidegree map is injective rationally:
\begin{equation}\label{eq:C3be-Pic-degree}
 \operatorname{Pic}(Z)_{\mathbb Q}\longrightarrow\mathbb Q^c,
 \qquad L\longmapsto(\deg(L|_{E_i}))_{i=1}^{c}.
\end{equation}
We have not inferred vanishing of \([L]-1\) merely from
torsion of \(L|_E\): we first took a power whose restriction
is actually trivial and then used the determinant of a
torsion \(K_0\)-class.

\emph{4. Exceptional divisors span the multidegrees.}
Let
\[
              M_{ij}=\deg\bigl(\mathcal O_Z(E_j)|_{E_i}\bigr).
\]
This matrix is symmetric; for \(i\ne j\) its entries are the
nonnegative local intersection lengths. If
\(a=(a_1,\ldots,a_c)^{\mathsf t}\) is supplied by
\eqref{eq:C3be-negative-ample}, every entry of \(Ma\) is
strictly negative. Put
\(Q=\operatorname{diag}(a)M\operatorname{diag}(a)\), and let
\(b_i=\sum_jQ_{ij}<0\). For \(x\in\mathbb R^c\),
\[
 x^{\mathsf t}Qx
       =\sum_i b_i x_i^2
          -\sum_{i<j}Q_{ij}(x_i-x_j)^2.
\]
Thus \(M\) is negative definite and invertible. The multidegrees
of \(\mathcal O_Z(E_j)\) span \(\mathbb Q^c\).
By \eqref{eq:C3be-Pic-degree}, these exceptional line bundles
generate \(\operatorname{Pic}(Z)_{\mathbb Q}\).

Restriction \(\operatorname{Pic}(Z)\to\operatorname{Pic}(U)\)
is surjective: represent a line bundle on the regular integral
scheme \(U\) by a Cartier divisor and take its Weil-divisor
closure on regular \(Z\), where it is Cartier.
Every exceptional line bundle restricts trivially. Hence
\begin{equation}\label{eq:C3be-surface-Pic-zero}
                     \operatorname{Pic}(U)_{\mathbb Q}=0.
\end{equation}
Chow localization identifies
\(\mathrm{CH}_1(\operatorname{Spec}B)\) with
\(\mathrm{CH}_1(U)=\operatorname{Pic}(U)\): the removed
closed point has \(\delta\)-dimension zero and \(U\) is regular.
This proves the normal-surface assertion for every excellent
normal local domain with residue field \(k\), including mixed
characteristic. No henselian hypothesis was used.

\emph{5. Every curve is killed on a normalized surface quotient.}
Return to the possibly nonnormal domain \(R\). For a curve
prime \(\mathfrak p\) with \(\dim(R/\mathfrak p)=1\),
catenarity supplies a prime
\[
              \mathfrak q\subset\mathfrak p,\qquad
                              \dim(R/\mathfrak q)=2.
\]
For \(N=2\) take \(\mathfrak q=(0)\); in general use the
prime immediately preceding \(\mathfrak p\) in a saturated
chain from \((0)\) to \(\mathfrak m\). Set \(B=R/\mathfrak q\)
and let \(T\) be its normalization. Excellence makes \(T\)
a finite normal integral \(B\)-algebra. It is semilocal, with
finitely many maximal ideals \(\mathfrak n_1,\ldots,\mathfrak n_s\);
each residue field is \(k\), since it is a finite extension
of the algebraically closed field \(k\).

For each \(i\), the ring \(T_{\mathfrak n_i}\) is excellent,
normal, local, and has dimension at most two. If its dimension
is two, Steps~1--4 give
\(\operatorname{Cl}(T_{\mathfrak n_i})_{\mathbb Q}=0\).
In dimension at most one it is a field or a discrete valuation
ring, and its class group is already zero.
For a Weil divisor \(G\) on \(\operatorname{Spec}T\), choose
a common positive integer \(e\) killing its classes at all
these maximal local rings. The rank-one coherent reflexive
module associated with \(eG\) is then free at every maximal
localization, and hence is an invertible module on
\(\operatorname{Spec}T\); we use the Weil-divisor/reflexive-module
description of \citet[Lemma 31.30.2, Tag 0EBM]
{frC3beStacksReflexiveClassGroup}. Every invertible module
over a semilocal ring is free by
\citet[Lemma 10.78.7, Tag 02M9]{frC3beStacksSemilocalPic}.
Indeed, the Chinese remainder theorem chooses a section
generating every residue fibre, and Nakayama's lemma makes it
a generator at every maximal ideal. Therefore \(eG\) is
principal and \(\operatorname{Cl}(T)_{\mathbb Q}=0\).

The dimension function induced from \(\operatorname{Spec}B\)
has value two at the generic point of \(\operatorname{Spec}T\):
the residue-field transcendence degrees of a finite map are
zero \citep[Lemma 29.53.3, Tag 02JW]
{frC3beStacksFiniteDimensionFunction}. Consequently
\begin{equation}\label{eq:C3be-semilocal-surface}
 \mathrm{CH}_1(\operatorname{Spec}T)_{\mathbb Q}
       =\operatorname{Cl}(T)_{\mathbb Q}=0,
\end{equation}
using the explicit dimension-function identification in
\citet[Section 42.24, immediately after Definition 42.24.1,
Tag 02SI]{frC3beStacksClassChow}.
This step needs neither a single branch of the normalization
nor any claim that all its maximal ideals have the same height.

Choose a prime \(\mathfrak P\subset T\) above
\(\mathfrak p/\mathfrak q\). It is neither zero nor maximal,
and hence has height one. For the finite morphism consisting
of normalization followed by the surface closed immersion,
\[
              g:\operatorname{Spec}T
                            \longrightarrow\operatorname{Spec}R,
\]
proper pushforward gives
\begin{equation}\label{eq:C3be-curve-push}
       g_*[V(\mathfrak P)]
          =[\kappa(\mathfrak P):\kappa(\mathfrak p)]
                                      [V(\mathfrak p)].
\end{equation}
The coefficient is a positive finite integer. Proper pushforward
respects rational equivalence in the dimension-function setting
\citep[Lemma 42.20.3, Tag 02S2]{frC3beStacksProperPush}.
The left side of \eqref{eq:C3be-curve-push} is zero rationally
by \eqref{eq:C3be-semilocal-surface}. Therefore every curve generator
of \(\mathrm{CH}_1(S)_{\mathbb Q}\) vanishes.

Finally, Chow localization
\citep[Lemma 42.19.3, Tag 02RX]{frC3beStacksChowLocalization}
gives \(\mathrm{CH}_1(S)\simeq\mathrm{CH}_1(X)\), because
the closed point carries no dimension-one cycles. On the
catenary local domain the codimension of a point of \(X\)
is unchanged on deleting \(\mathfrak m\), so this is exactly
\(\mathrm{CH}^{N-1}(X)\). This proves
\eqref{eq:C3be-top-vanishing} and any asserted quotient
vanishing. We claim neither integral vanishing nor a uniform
annihilating integer.
\end{proof}

The finite-type equal-characteristic surface case has established
antecedents: \citet[Theorem 4.5, p. 42, proof and Remark 4.9,
p. 43]{frC3beTanaka2014} proves \(\mathbb Q\)-factoriality of
normal surfaces over \(\overline{\mathbb F}_p\), and the proof
of \citet[Theorem 2.4, author pp. 4--5]{frC3beEsser2023}
proves torsion of their henselian local class groups, crediting
Artin's argument. We give the determinant and surface-quotient
proof above to justify the stated arbitrary-excellent and
mixed-characteristic extension; we make no priority claim.
